\documentclass[../../../include/open-logic-section]{subfiles}

\begin{document}

\olfileid[hi]{sth}{choice}{tarskiscott}
\olsection{टार्स्की--स्कॉट युक्ति}

\olref[cardinals][cardsasords]{defcardinalasordinal} में हमने कार्डिनलों
को क्रमसूचकों के रूप में परिभाषित किया। इसके लिए हमने सुव्यवस्था अभिगृहीत
माना। हमने ऐसा केवल इसलिए किया कि यही ``प्रचलित मानक'' है।

इसलिए सुव्यवस्था से जुड़े दार्शनिक प्रश्नों पर चर्चा करने से पहले यह स्पष्ट
होना महत्त्वपूर्ण है कि हम प्रचलित मानक से हट \emph{सकते हैं} और सुव्यवस्था
माने \emph{बिना} कार्डिनलों का सिद्धांत विकसित कर सकते हैं। हम फिर भी
$\cardeq{A}{B}$, $\cardle{A}{B}$ और $\cardless{A}{B}$ की परिभाषाएँ वैसे
ही उपयोग कर सकते हैं जैसे वे \olref[sfr][siz][]{chap} में आई थीं। हमें केवल
\emph{कार्डिनल} की एक नई धारणा चाहिए।

एक सहज विचार $A$ की कार्डिनलता को इस प्रकार परिभाषित करने का प्रयास होगा:
\[
	\Setabs{x}{\cardeq{A}{x}}.
\]
आप इसकी तुलना फ़्रेगे की $\#xFx$ की परिभाषा से कर सकते हैं, जिसकी रूपरेखा
\olref[cardinals][hp]{sec} के बिल्कुल अंत में दी गई थी। वहाँ संकेत किए कारणों
से यह परिभाषा विफल होती है। प्रत्येक एकल समुच्चय $\{\emptyset\}$ के साथ
समसंख्य है। किंतु पदानुक्रम के प्रत्येक उत्तरवर्ती चरण पर नए एकल समुच्चय
बनते हैं (बस पिछले चरण के एकल समुच्चय पर विचार करें)। अतः
$\Setabs{x}{\cardeq{A}{x}}$ का अस्तित्व नहीं है, क्योंकि उसका कोई रैंक नहीं
हो सकता।

इस समस्या से बचने के लिए हम टार्स्की और स्कॉट की एक युक्ति का उपयोग करते
हैं:\footnote{स्मरण रहे: सभी सूत्रों में प्राचल हो सकते हैं (जब तक स्पष्ट
रूप से अन्यथा न कहा गया हो)।}

\begin{defn}[Tarski--Scott]
किसी भी सूत्र $\phi(x)$ के लिए $[x:\phi(x)]$ उन सभी लघुतम संभव रैंक वाले
$x$ का समुच्चय हो जिनके लिए $\phi(x)$ (अथवा यदि कोई ऐसा $x$ न हो तो
$\emptyset$)।
\end{defn}

हमें जाँचना चाहिए कि यह परिभाषा वैध है। $\ZF$ में काम करते हुए
\olref[spine][foundation]{zfentailsregularity} प्रत्याभूत करता है कि
प्रत्येक $x$ के लिए $\setrank{x}$ का अस्तित्व है। अब यदि $\phi$ को संतुष्ट
करने वाली कोई सत्ताएँ हैं, तो $\alpha$ वह लघुतम रैंक हो सकता है जिसके लिए
$(\exists x\subseteq V_\alpha)\phi(x)$, अर्थात्
$(\forall\beta\in\alpha)(\forall x\subseteq V_\beta)\lnot\phi(x)$। तब हम
पृथक्करण द्वारा $[x:\phi(x)]$ को
$\Setabs{x\in V_{\alpha+1}}{\phi(x)}$ के रूप में परिभाषित कर सकते हैं।

टार्स्की--स्कॉट युक्ति का औचित्य देने के बाद अब हम इसका उपयोग कार्डिनलता की
एक धारणा परिभाषित करने में कर सकते हैं:

\begin{defn}
$A$ की \textsc{ts}-कार्डिनलता
$\text{tsc}(A)=[x:\cardeq{A}{x}]$ है।
\end{defn}

\textsc{ts}-कार्डिनल की परिभाषा सुव्यवस्था का उपयोग नहीं करती। किंतु उस
अभिगृहीत के बिना भी हम दिखा सकते हैं कि \emph{\textsc{ts}-कार्डिनल}
\olref[cardinals][cardsasords]{defcardinalasordinal} में परिभाषित
\emph{कार्डिनलों} की तरह काफ़ी व्यवहार करते हैं। उदाहरण के लिए यदि हम
\olref[cardinals][cardsasords]{lem:CardinalsBehaveRight} और
\olref[card-arithmetic][opps]{lem:SizePowerset2Exp} को
\textsc{ts}-कार्डिनलों के पदों में फिर से कहें, तो प्रमाण सुव्यवस्था माने
बिना $\ZF$ में ठीक चलते हैं।

इसी विषय में यह ध्यान देने योग्य है कि टार्स्की--स्कॉट युक्ति से हम
क्रमसूचकों का सिद्धांत भी विकसित कर सकते हैं। जहाँ $\tuple{A,<}$ कोई
सुव्यवस्था है, वहाँ
$\text{tso}(A,<)=[\tuple{X,R}:\ordeq{\tuple{A,<}}{\tuple{X,R}}]$ रखें।
कार्डिनलों और क्रमसूचकों के इस निरूपण पर अधिक के लिए
\citet[chs.~9--12]{Potter2004} देखें।

\end{document}
